As an exercise of image processing, this problem involves use of a simple
calculator and tabular data (table 1.1) which contains the pixel values of an
image during the given exposure time. This image, which is a part of a larger
CCD image, was taken by a small CCD camera, installed on an amateur telescope
and using a $V$ band filter. Figure 1.1 shows this $50 \times 50$ pixels
image that contains 5 stars.

\ioaafig{2009_DA_D1-f1.pdf}

% TABLE ABRIDGED — table 1.1 (50x50 grid of pixel values) in crop edition

In table 1.1 the first row and column indicates the pixels' x and y
coordinates. Table 1.2 gives the telescope and the image specifications.

\begin{center}
Table 1.2\\[2pt]
\begin{tabular}{|l|c|}
\hline
Telescope focal length & $1.20$\,m \\
\hline
CCD pixel size & $25 \times 25\,\mu$m \\
\hline
Exposure time & $450$\,s \\
\hline
Telescope zenith angle & $25^\circ$ \\
\hline
Average extinction coefficient in $V$ band & 0.3\,mag/airmass \\
\hline
\end{tabular}
\end{center}

The observer identified stars 1, 3 and 4 by comparing this image with
standard star catalogues. Table 1.3 shows stars true magnitudes ($m_t$) as
given in the catalogue.

\begin{center}
Table 1.3\\[2pt]
\begin{tabular}{|c|c|}
\hline
Star & $m_t$ \\
\hline
1 & 9.03 \\
\hline
3 & 6.22 \\
\hline
4 & 8.02 \\
\hline
\end{tabular}
\end{center}

\begin{description}
\item[(a)] Using the available data, determine the instrumental magnitudes
  of the stars in the image. Assume the dark current is negligible and the
  image is flat fielded. For simplicity you can use a square aperture.

  Hint: The instrumental magnitude is calculated using the difference
  between the measured flux from the star in the aperture and the flux from
  an equivalent area of dark sky.
\item[(b)] The instrumental magnitude of a star in a CCD image is related to
  true magnitude as
  \[ m_I = m_t + KX - Zmag \]
  where $K$ is extinction, $X$ is airmass, $m_I$ and $m_t$ are respectively
  instrumental and true magnitude of star and $Zmag$ is zero point constant.
  Calculate the zero point constant ($Zmag$) for identified stars. Calculate
  average zero point constant ($\overline{Zmag}$).

  Hint: Zero point constant is the constant reducing extinction-free
  magnitudes to the true magnitude.
\item[(c)] Calculate true magnitudes of stars 2 and 5.
\item[(d)] Calculate CCD pixel scale for the CCD camera in units of arcsec.
\item[(e)] Calculate average brightness of dark sky in magnitude per square
  arcsec ($m_{sky}$).
\item[(f)] Use a suitable plot to estimate astronomical seeing in arcsec.
\end{description}
